\subsection{Cold start i cachowanie}

Pomiar czasu wykonania zapytania w nowoczesnym SZBD jest zadaniem
nietrywialnym, ponieważ silnik bazodanowy jest systemem wielowarstwowym,
w którym powtórne wykonanie identycznego zapytania może być wielokrotnie
szybsze od pierwszego -- a różnica wynika nie z optymalizacji samej
operacji, lecz ze stanu cache'ów na różnych warstwach abstrakcji. W bazach
testowanych w niniejszej pracy występują co najmniej cztery niezależne
warstwy buforowania: bufor stron silnika (\texttt{shared\_buffers}
PostgreSQL, \texttt{innodb\_buffer\_pool} MySQL, \texttt{wiredTigerCache}
MongoDB, \texttt{pagecache} Neo4j), cache planów zapytań (\texttt{plan\_cache}
PostgreSQL, \texttt{query plan cache} MongoDB, \texttt{Cypher query cache}
Neo4j), bufor rezultatów po stronie drivera oraz cache stron systemu
operacyjnego utrzymywany przez jądro Linuksa kontenera Docker. Każda
z tych warstw ma inny czas życia i inną politykę unieważniania.

Aby pomiary były reprezentatywne dla typowego cyklu pracy aplikacji
(scenariusz wykonywany wielokrotnie z różnymi parametrami, na bazie
,,rozgrzanej'' przez wcześniejsze obciążenie), przyjęto następującą
dyscyplinę pomiarową. Każdy scenariusz uruchamiany jest w
\textbf{trzech następujących po sobie próbach} (\texttt{trials}=3),
a do dalszej analizy przekazywane są wszystkie trzy wyniki -- mediana
służy do agregacji w wizualizacjach, a pojedyncze trafienia są
zachowywane w pliku CSV razem z numerem próby. Zastosowanie mediany
zamiast średniej arytmetycznej jest świadomą decyzją -- ogranicza
wpływ pojedynczych odchyleń wynikających z chwilowych obciążeń
systemu operacyjnego (np. zrzutu dziennika WAL na dysk, krótkiej
pracy procesu w tle), które średnia by przesuwała.

\subsubsection*{Czyszczenie cache zapytań przed scenariuszem}

Przed pierwszą próbą każdego scenariusza, dla każdej z czterech baz,
wykonywana jest dedykowana procedura \texttt{\_flush\_caches()}
(\texttt{src/benchmarks/runner.py:186}) czyszcząca cache planu zapytań
silnika -- bez zerowania bufora stron. Cel jest precyzyjny: chcemy
wyeliminować wpływ planu zapamiętanego z poprzedniego scenariusza
(np.~zapytanie INSERT mogłoby skompilować plan, który dla zapytania
SELECT byłby nieoptymalny), zachowując jednocześnie ,,rozgrzaną''
warstwę pamięci -- co oddaje rzeczywiste warunki produkcyjne, w
których baza nie jest restartowana między zapytaniami klienta.
Konkretne instrukcje dla każdego silnika:

\begin{itemize}
    \item \textbf{PostgreSQL} -- \texttt{DISCARD ALL} -- oficjalna
    instrukcja zwalniająca cache zapytań przygotowanych, kursory
    sesji, zmienne sesyjne i tymczasowe tabele dla bieżącej sesji.
    Wykonywana w trybie \texttt{autocommit=True}, ponieważ w
    transakcji jawnej nie jest dopuszczalna.
    \item \textbf{MySQL} -- \texttt{FLUSH TABLES} -- zamyka wszystkie
    otwarte tabele i unieważnia ich metadane w cache silnika.
    \item \textbf{MongoDB} -- polecenie \texttt{planCacheClear}
    wykonywane na każdej z głównych kolekcji
    (\texttt{users}, \texttt{content}, \texttt{watch\_history},
    \texttt{ratings}, \texttt{payments}). Mongo per-kolekcja
    utrzymuje cache wybranego planu wykonania -- bez wyczyszczenia
    drugi i trzeci pomiar mogłyby korzystać z innego planu niż pierwszy.
    \item \textbf{Neo4j} -- \texttt{CALL db.clearQueryCaches()} -- procedura
    zerująca cache skompilowanych zapytań Cypher na poziomie całej bazy.
\end{itemize}

\noindent Pierwszy z trzech \texttt{trials} każdego scenariusza pełni
więc funkcję \emph{cold start} po stronie kompilacji planu, ale jest
ciepły po stronie bufora stron. Próby drugą i trzecią należy traktować
jako reprezentujące zapytania w bazie ,,utrzymywanej w cyklu''.
Rozkład wartości \texttt{trial=1}~vs.~\texttt{trial=2,3} sam w sobie
niesie informację o stabilności silnika -- w scenariuszach z dużą
różnicą widać, że cache planów był kluczowy.

\subsubsection*{Świadomie nieczyszczone warstwy}

Bufor stron silnika (\texttt{shared\_buffers}, \texttt{buffer\_pool},
\texttt{wiredTigerCache}, \texttt{pagecache}) oraz cache stron systemu
operacyjnego pozostają nietknięte. Decyzja jest celowa -- realny system
produkcyjny nie restartuje się przed każdym zapytaniem, a wykres
,,zimny vs ciepły'' opisany w dziedzinie pamięci stronowej jest
ortogonalny do pytania badawczego pracy (wpływ indeksów na czas
zapytań). Drugorzędną zaletą jest to, że pomiary dla wolumenów
\emph{medium} i \emph{large} mieszczą się w realistycznym czasie --
wymuszanie cold-startu pamięci stronowej dla 10~mln wierszy w
każdym z 24~scenariuszy wydłużyłoby pojedynczy run benchmarka o
rzędy wielkości i nie zmieniłoby relacji między bazami.

\subsubsection*{Pomiar czasu}

Pomiar czasu wykonuje funkcja \texttt{time.perf\_counter\_ns()}
biblioteki standardowej Pythona, zwracająca licznik nanosekundowy
o najwyższej dostępnej rozdzielczości monotonicznej zegara systemowego.
Wynik jest dzielony przez \(10^{6}\) i zapisywany jako liczba
zmiennoprzecinkowa milisekund z trzema miejscami po przecinku.
Mierzony jest \emph{wyłącznie} blok \texttt{run()} scenariusza --
przygotowanie danych w \texttt{setup()} oraz porządki w \texttt{teardown()}
są wykonywane poza pomiarem, dzięki czemu raportowany czas obejmuje
tylko właściwą operację CRUD razem z jej narzutem komunikacyjnym
(serializacja parametrów, RTT do silnika, parsowanie/optymalizacja
po stronie silnika, wykonanie, deserializacja wyniku).

% -------------------------------------------------------
\subsection{Pipeline testowy}

Cały cykl badawczy został scalony w jedno polecenie CLI --
\texttt{python main.py run-all --volume \{small|medium|large\}} --
które wykonuje 8-krokowy pipeline rozdzielający fazę bez indeksów
od fazy z indeksami:

\begin{enumerate}
    \item \texttt{generate} -- wygenerowanie plików CSV
    z deterministycznym ziarnem 42.
    \item \texttt{load --no-indexes} -- załadowanie wszystkich czterech
    baz, bez tworzenia indeksów wtórnych.
    \item \texttt{benchmark} -- 24~scenariusze $\times$ 4~bazy
    $\times$ 3~próby na bazie bez indeksów.
    \item \texttt{explain} -- analiza planów wykonania scenariuszy
    SELECT (rozdział 10) dla wariantu bez indeksów.
    \item \texttt{load} -- ponowne załadowanie tych samych danych,
    tym razem ze stworzeniem indeksów wtórnych.
    \item \texttt{benchmark --with-indexes} -- powtórzenie tych
    samych 24~scenariuszy w wariancie z indeksami.
    \item \texttt{explain --with-indexes} -- analiza planów
    wykonania w wariancie z indeksami.
    \item \texttt{visualize} -- agregacja plików CSV z wynikami
    i generacja kompletu wykresów PNG (heatmapy, słupkowe per
    kategoria CRUD, wykresy skalowalności).
\end{enumerate}

Krok 5 -- ponowne załadowanie danych -- jest istotny merytorycznie:
zamiast budować indeksy na bazie, którą już ,,obciążyły'' pomiary
fazy bez indeksów (z fragmentacją tabel po DELETE, ze zmienionymi
wartościami sekwencji po INSERT, z zaktualizowanymi statystykami
analizatora), pipeline odtwarza świeży, deterministyczny stan, na
którym dopiero wykonuje benchmarki z indeksami. Dzięki temu pomiary
\texttt{not\_indexed} i \texttt{with\_indexes} mierzą \emph{ten sam}
zbiór danych -- różni je tylko obecność indeksów.

\subsubsection*{Architektura modułu \texttt{BenchmarkRunner}}

Sercem warstwy pomiarowej jest klasa \texttt{BenchmarkRunner}
(\texttt{src/benchmarks/runner.py}). Po inicjalizacji parametrami
(\texttt{volume}, \texttt{results\_dir}) jej cykl życia obejmuje
trzy fazy:

\begin{enumerate}
    \item \texttt{connect()} -- nawiązanie połączeń z czterema
    silnikami (psycopg, PyMySQL, PyMongo, neo4j-driver) na
    \texttt{localhost} z loginami i hasłami zdefiniowanymi w
    \texttt{docker-compose.yml}. Połączenia są zachowywane na czas
    całego runa.
    \item \texttt{build\_context()} -- zbudowanie obiektu
    \texttt{BenchmarkContext} przechowującego maksymalne wartości
    kluczy głównych każdej tabeli (odczytane przez
    \texttt{SELECT MAX(id) FROM \dots} z PostgreSQL, używanego jako
    referencyjne źródło). Identyfikatory te są niezbędne, żeby
    scenariusze \texttt{INSERT} generowały klucze wykraczające
    poza istniejący zbiór (\texttt{ctx.test\_id(table) =
    max\_ids[table] + 100\,000 + counter}), a \texttt{UPDATE}
    i \texttt{DELETE} mogły operować na losowych, faktycznie
    istniejących wierszach (\texttt{ctx.random\_id(table)}).
    \texttt{BenchmarkContext} dzielony jest między wszystkie
    scenariusze całego runa.
    \item \texttt{run\_all(with\_indexes, trials, scenario\_ids)}
    -- pętla pomiarowa.
\end{enumerate}

Wewnątrz \texttt{run\_all()} pętle są zagnieżdżone w kolejności:
\textbf{scenariusz~$\rightarrow$~baza~$\rightarrow$~próba}.
Dla każdej pary (scenariusz, baza) tworzone jest nowe, świeże
połączenie metodą \texttt{\_create\_connection()} (zamykane po
zakończeniu trzech prób), a sam obiekt scenariusza jest płytko
kopiowany (\texttt{copy.copy(scenario)}) -- co izoluje stan
wewnętrzny scenariusza w zakresie pojedynczego pomiaru.
Dopiero po stworzeniu połączenia i kopii scenariusza wywoływana
jest procedura \texttt{\_flush\_caches()} opisana w sekcji 9.1,
po czym następują trzy iteracje pętli prób.

Każda próba ma identyczną strukturę zgodną z protokołem klasy
bazowej \texttt{BaseScenario}:

\begin{lstlisting}[language=Python,basicstyle=\ttfamily\small]
sc.setup(db_name, conn, ctx)        # poza pomiarem
start = time.perf_counter_ns()
sc.run(db_name, conn, ctx)          # MIERZONY blok
elapsed_ms = (time.perf_counter_ns() - start) / 1_000_000
sc.teardown(db_name, conn, ctx)     # poza pomiarem
\end{lstlisting}

\noindent \texttt{setup()} przygotowuje stan bazy potrzebny do
operacji (np. dla scenariusza UPDATE wybiera identyfikator
istniejącego wiersza), \texttt{run()} wykonuje właściwą operację
CRUD i jest jedynym mierzonym blokiem, a \texttt{teardown()}
przywraca bazę do stanu wyjściowego (np.~rolluje INSERT)
-- aby trzy próby tego samego scenariusza były porównywalne i
nie kumulowały efektów ubocznych. Każda z trzech metod
posiada specjalizację per-silnik
(\texttt{run\_postgres}, \texttt{run\_mysql}, \texttt{run\_mongo},
\texttt{run\_neo4j}); klasa bazowa rozdziela wywołania do
odpowiedniej implementacji przez \texttt{getattr(self,
f"run\_\{db\_type\}")}.

\subsubsection*{Obsługa błędów}

Każda z faz \texttt{setup}/\texttt{run}/\texttt{teardown} jest opakowana
osobnym blokiem \texttt{try-except}. Wystąpienie wyjątku w którejkolwiek
z nich powoduje wypisanie pełnego stack trace, próbę
\texttt{conn.rollback()} dla baz transakcyjnych
(\texttt{\_try\_rollback()} dla PostgreSQL i MySQL) oraz pominięcie
zapisu wyniku tej konkretnej próby. Pętla nie przerywa się -- pozostałe
próby i scenariusze kontynuują wykonanie. Dzięki temu pojedynczy błąd
w scenariuszu (np. zerwanie połączenia, deadlock) nie unieważnia
całego runa, tylko ogranicza liczbę wartości w danej komórce wyników.

\subsubsection*{Format wyników}

Wynikiem każdego runa benchmarka jest plik
\texttt{benchmark\_\{volume\}\_\{label\}.csv} (gdzie
\texttt{label}~$\in$~\{\texttt{no\_indexes}, \texttt{with\_indexes}\})
zapisany w katalogu \texttt{results/}. Każdy wiersz pliku to jedna
próba pojedynczego scenariusza na pojedynczej bazie -- nie agregat.
Zestaw kolumn:

\begin{itemize}
    \item \texttt{scenario\_id} -- identyfikator scenariusza
    (np. \texttt{S1}, \texttt{I7}, \texttt{U3}, \texttt{D2}),
    \item \texttt{scenario\_name} -- czytelna nazwa,
    \item \texttt{category} -- jedna z czterech kategorii
    (\texttt{INSERT}, \texttt{SELECT}, \texttt{UPDATE},
    \texttt{DELETE}),
    \item \texttt{database} -- \texttt{postgres},
    \texttt{mysql}, \texttt{mongo}, \texttt{neo4j},
    \item \texttt{volume} -- \texttt{small} / \texttt{medium}
    / \texttt{large},
    \item \texttt{trial} -- numer próby (1--3),
    \item \texttt{time\_ms} -- zmierzony czas
    \texttt{run()} w milisekundach,
    \item \texttt{with\_indexes} -- flaga logiczna
    \texttt{True}/\texttt{False}.
\end{itemize}

\noindent Dla pełnego pipeline'u na pojedynczym wolumenie pojedynczy
plik z wynikami zawiera $24 \times 4 \times 3 = 288$~wierszy, czyli łącznie
576 pomiarów (\emph{not\_indexed} + \emph{with\_indexes}).
Sumarycznie pełny przebieg na trzech wolumenach generuje
$3 \times 576 = 1\,728$~rekordów wyników, z których wszystkie są
przetwarzane wspólnie przez moduł \texttt{Visualizer}.

\subsubsection*{Agregacja wyników}

Moduł \texttt{src/analysis/visualizer.py} wczytuje pliki CSV za
pomocą \texttt{pandas.read\_csv} i wykonuje agregację
\texttt{groupby} po polach \texttt{(scenario\_id, scenario\_name,
category, database, volume, with\_indexes)}. Funkcja agregująca to
\texttt{median} po kolumnie \texttt{time\_ms}, uzupełniona przez
\texttt{min} i \texttt{max} z trzech prób -- co pozwala każdemu
wykresowi prezentować zarówno wartość typową, jak i rozrzut.
Mediana z trzech wartości jest oczywiście podatna na uproszczenie
(jest po prostu wartością środkową), ale -- biorąc pod uwagę
wcześniej opisaną dyscyplinę cold-startu cache planów -- to
wartość, którą najtrudniej zafałszować pojedynczemu wahaniu.

W rozdziałach 11~i~12 raportowane są w tej kolejności: mediana
trzech prób (jako wartość punktowa wykresu), a następnie zakres
\texttt{min--max} (jako wąsy lub jasny zakres pod krzywą),
a w rozdziale poświęconym weryfikacji hipotezy -- stosunek
median wariantu \texttt{with\_indexes} do mediany wariantu
\texttt{not\_indexed} dla każdej pary (scenariusz, baza), zwizualizowany
na mapach cieplnych.
